\section{Caps and a nonzero ordinary pairing}\label{sec:caps}

This section constructs a nonzero ordinary pairing and puts every
Donaldson insertion in two fixed end caps.  These caps supply the fixed
states paired with the repeated cobordism maps in Section~\ref{sec:stack}.
The intervening cobordisms will carry no moving surface insertions.
We retain the determinant-one gauge group and the fixed collar
identifications of Section~\ref{sec:ordinary}.  Write $x_0$ for the point
class and use the normalization
\begin{equation}\label{eq:caps:normalization}
 \mu(\xi)=-\tfrac14 p_1(\ad E)/\xi .
\end{equation}
A \emph{tight surface} is an embedded oriented surface $B$ with
$B^2=2g(B)-2>0$.  The capped minimal Seifert surface $A\subset Y_0$ has
genus two and square zero in the closures below.

\begin{theorem}[Fixed caps with a nonzero pairing]\label{thm:caps:pairing}
There are two fixed compact caps $C_-$ and $C_+$ with common rational
boundary $Y_2$, fixed determinant data trivialized at that boundary, and
finite collections of integral point and surface probes in their interiors
with the following properties.
\begin{enumerate}[label=\textup{(\roman*)}]
\item Each cap has $H_1(C_\pm;\Z)=0$ and $b^+(C_\pm)>0$.  Each contains
a square-zero copy of $A$ on which the determinant class evaluates as $1$,
and a tight genus-two surface of square two.
\item Each cap has one additional, separate exceptional summand.  The
determinant evaluates as zero on its exceptional sphere, and every probe
is supported in the old part of the cap.
\item For suitable generic cap data and small ordinary Floer data at the
rational boundary, the cut counts give a rational cycle $x$ and a rational
cocycle $y$ on the irreducible determinant-one Floer complex, with
\begin{equation}\label{eq:caps:nonzero-pairing}
                         yx\ne0.
\end{equation}
Fixed integral multiples of these are an integral cycle and cocycle.
\item The total real degree $z$ of all cap probes is fixed.  Reducible
cap limits in any of the bounded charge ranges used below are excluded by
the cap period choices.  The period choices admit the finite isolation
extensions of Proposition~\ref{prop:estimates:cap}.
\end{enumerate}
The cycle and cocycle can be recomputed with the final cap data used in
Section~\ref{sec:stack}; their pairing is the same nonzero closed
Donaldson count, up to the fixed normalization and orientation sign.
\end{theorem}

We first construct a closure, prove that it has a nonzero Donaldson
invariant, and move its insertions away from a rational seam.  The exact
sections used to represent the insertions are supplied by
Proposition~\ref{prop:analysis:ordinary-library}; the present section proves
the labelled identities, the compactness budgets, and the comparison with
closed Donaldson invariants.  This use of the ordinary analytic library
does not require the spinor compactification or the nonzero count to be
constructed later.

\subsection{The closure and its tight surfaces}

The foliation, contact, and symplectic-cap route to a closed
nonvanishing problem is assembled in
Kronheimer--Mrowka~\cite[Theorem~8, Lemmas~9--14, and
Proposition~15]{KronheimerMrowka2004}.  The construction below establishes
simple connectivity and the tight surfaces used in the nonvanishing
argument.  It also gives the stated geography bounds.

\begin{lemma}[Enhanced symplectic closure]\label{lem:caps:closure}
There is a symplectic closure $X_0$ of $Y_0$ such that its two halves are
simply connected, $A$ is symplectic, and each half contains a tight
genus-two surface $B$.  The closure can be chosen to satisfy
\begin{equation}\label{eq:caps:geography}
 q(X_0):=\frac{\chi(X_0)+\sigma(X_0)}4\ge3,
 \qquad s(X_0):=2\chi(X_0)+3\sigma(X_0)\ge q(X_0)-2.
\end{equation}
Moreover, $[A]$ is primitive and there is $w\in H^2(X_0;\Z)$ with
$\langle w,A\rangle=1$.
\end{lemma}

\begin{proof}
Gabai's taut-foliation construction for zero surgery on a nontrivial knot
allows the capped norm-minimizing Seifert surface to be a leaf
\cite[Corollary~8.10 and Remark~8.13]{Gabai1987}.
The contact approximation and weak symplectic product
construction preserve positivity on this leaf
\cite{EliashbergThurston1998}.  Thus a product neighborhood of $Y_0$
embeds in a symplectic cobordism, with $A$ symplectic and outer boundaries
weakly compatible with their contact structures.  We describe the additional
cap choices; they strengthen
the homological conclusions alone of the usual cap construction
\cite[Theorem~8 and Lemmas~9--12]{KronheimerMrowka2004}.

Choose a supporting open book with connected binding at each end
\cite[Theorems~3--4]{Giroux2002}.
Positive stabilizations can increase the page genus while restoring
connected binding in pairs of stabilizations.  We may therefore take its
genus to be
\[
                     G=\frac{(d-1)(d-2)}2
\]
for an arbitrarily large $d\ge6$.  Attach Eliashberg's symplectic binding
handle \cite[Theorem~1.1]{Eliashberg2004}.  The new boundary is a bundle of
closed genus-$G$ surfaces over the circle, with a prescribed closed
boundary two-form positive on the fibers.  Viewed from this new boundary,
the handle cobordism has only a relative two-handle.  In particular its
fundamental group is generated by the image of that boundary.

Factor the required inverse boundary monodromy into positive twists on
nonseparating curves~\cite[Lemma~2, arXiv version~2]{AkbulutOzbagci2002}.
Such a factorization exists for the closed fiber: nonseparating twists
generate the closed surface mapping class group, and a positive chain
relation expressing the identity contains each of the standard generators.
Solving this relation for a chosen inverse expresses that inverse as a
positive word; conjugation does the same for any nonseparating twist.
Append as many positive identity factorizations as necessary from generic
degree-$d$ plane pencils.  These pencils have only irreducible nodal
singular fibers.  Their blown-up total spaces are
$\CP^2\#d^2\overline{\CP}^{\,2}$ and have sections.  Removing a regular
fiber leaves a simply connected disk fibration: its fundamental group is
normally generated by the meridian of the removed fiber, and a punctured
section is a disk bounding this meridian.  Equivalently, the vanishing
cycles of one plane-pencil identity block normally generate the fiber
fundamental group.  Appending one such block consequently kills the
fundamental group of the disk fibration with the originally prescribed
monodromy.  Let $Z$ denote the resulting simply connected disk filling.

There is no additional boundary-flux hypothesis in this construction.
Indeed, Poincar\'e--Lefschetz duality gives
\[
 H^3(Z,\partial Z;\R)\cong H_1(Z;\R)=0,
\]
so the class of the prescribed boundary two-form extends over $Z$.
It is positive on the regular fiber and on each singular fiber component,
since every singular fiber is irreducible.  Here is the relative form
construction with this particular boundary condition.  On a boundary
collar extend the prescribed two-form, and choose representatives of the
extended class positive on the vertical tangent planes in fiber
neighborhoods.  At a critical point use the usual local Lefschetz model
and a form taming that model.  Representatives in overlapping base
neighborhoods differ by exact forms after their fiber integrals have been
matched.  Patching their primitives with a partition of unity on the base
therefore preserves both closedness and vertical positivity; the correction
terms contain a base differential and vanish on vertical pairs.  Carry
out this patching relative to the boundary collar.  Adding a sufficiently
large pullback of an area form on the disk gives a symplectic form,
including on the finitely many critical neighborhoods.  That added form
restricts to zero on the tangent bundle of the boundary.  Thus its
tangential boundary restriction is exactly the prescribed one.  Equal
restrictions give compatible collars by the symplectic hypersurface
neighborhood construction.  This also agrees with the boundary-form and
flux arrangements in
\cite[Lemmas~3.4--3.5 and Remark~3.6]{Eliashberg2004}.

Gluing $Z$ kills the group of the binding-handle cobordism, because its
new boundary surjects on that group.  It follows that the resulting half
of the closure is simply connected.  Do this at both ends.  Van Kampen's
theorem then gives $\pi_1(X_0)=0$.

For completeness, two adjacent plane-pencil identity blocks give the
required tight surface inside either cap.  Conjugate their fiber
identifications so that a chosen nonseparating vanishing cycle agrees in
the two blocks, and put these chosen critical points next to the
intervening product region.  Join their thimbles across that region.  The
Lefschetz framings give an embedded sphere $S$ of square $-2$.  A curve in
the fiber dual to the vanishing cycle gives a rim torus $R$ in the product
region, with $R^2=0$ and $S\cdot R=1$.  Take disjoint parallel copies
$R_1,R_2$ and resolve their two intersections with $S$.  The resulting
connected embedded oriented surface satisfies
\[
 \chi(B)=2+0+0-4=-2,
 \qquad B^2=(S+R_1+R_2)^2=-2+4=2.
\]
Thus $g(B)=2$ and $B$ is tight.  This is the sphere-and-two-tori
construction of \cite[Corollary~8.5]{KronheimerMrowka1995}, in the
matching-thimble form used in
\cite[Lemma~4.6, arXiv version~2]{Sivek2015}.  It requires neither a new
cap theorem nor a genus-two nonvanishing theorem.

One additional plane-pencil block has Euler characteristic $3+d^2$
and signature $1-d^2$.  Fiber sum along genus $G$ adds another $4G-4$
to the Euler characteristic.  Consequently its increments are
\begin{equation}\label{eq:caps:block-increments}
                     \Delta q=G,
 \qquad \Delta s=1-d^2+8G.
\end{equation}
For $d\ge6$, $\Delta s-\Delta q=1-d^2+7G>0$.  Sufficiently many
blocks give \eqref{eq:caps:geography}.  In particular
$b^+(X_0)=2q(X_0)-1\ge5$.

Finally symplectic adjunction gives $K_\omega\cdot A=2$.  The second
homology of the simply connected $X_0$ is torsion free.  If $[A]=ka$
with $k>1$, this equality forces $k=2$ and $K_\omega\cdot a=1$.
But $a^2=0$, contradicting the fact that $K_\omega$ is characteristic.
Hence $[A]$ is primitive.  Unimodularity of the intersection form supplies
an integral $w$ with $\langle w,A\rangle=1$.
\end{proof}

\subsection{An extremal Floer line and Donaldson nonvanishing}

\begin{proposition}[Donaldson nonvanishing for the closure]\label{prop:caps:nonvanishing}
For the integral determinant $w$ of Lemma~\ref{lem:caps:closure},
there are $m,a\ge0$ and $h\in H_2(X_0;\Z)$ such that
\begin{equation}\label{eq:caps:initial-invariant}
       h\cdot A=1,
       \qquad D^w_{X_0}(x_0^m h^a)\ne0.
\end{equation}
\end{proposition}

We prove this by comparing positive Lefschetz words.  A pencil on a
blowup of $X_0$ supplies the word to be studied; an elliptic surface
supplies a model word with a nonzero Donaldson invariant.  On a
one-dimensional generalized Floer summand the model acts by a nonzero
scalar.  Conjugation and Hurwitz moves extract the pencil word from a
product of model words, forcing its action on that same line to be
nonzero.  Simple type and the blowup formula then recover the prescribed
determinant $w$ on $X_0$.

The argument uses ordinary irreducible gluing in
Lemma~\ref{lem:caps:irreducible-gluing}, the Donaldson structure theorem,
and the calculation for elliptic surfaces.  It does not identify the
Donaldson series of an arbitrary symplectic manifold with its
Seiberg--Witten series.  We begin with the generalized Floer summand
on which the comparison will take place.

\begin{lemma}[The product extremal generalized spaces]\label{lem:caps:extreme}
Let $V_g$ be the full determinant-one instanton homology over $\C$ of
$\Sigma_g\times S^1$, with determinant dual to $\{p\}\times S^1$,
where $g\ge2$.  Set $S=\mu(\Sigma_g)$ and $P=\mu(x_0)$.
Every joint eigenvalue $(\lambda,\nu)$ has $|\lambda|\le2g-2$.
Those with equality are exactly
\begin{equation}\label{eq:caps:extreme-spectrum}
       (\lambda,\nu)=\bigl(i^r(2g-2),(-1)^r2\bigr),
                  \qquad r=0,1,2,3.
\end{equation}
Each corresponding simultaneous generalized eigenspace has dimension one.
\end{lemma}

\begin{proof}
The gauge convention matters here.  The full group is the determinant-one
group of \cite[Section~7.1]{KronheimerMrowka2010}; Mu\~noz's ring
calculation uses the additional involution quotient.  The spectrum and the
conversion between these conventions are described in
\cite[Proposition~7.1 and Section~7.2]{KronheimerMrowka2010}.
Proposition~7.4 there states an ordinary eigenspace assertion.  We give the
extra local-algebra argument needed for a generalized eigenspace,
using the localization framework of
Mu\~noz~\cite[Section~5, equations~(6)--(7)]{Munoz1999Ring}.

In the quotient ring use Mu\~noz's coordinates
$\alpha=2S$, $\beta=-4P$, and $\gamma$.  Its invariant summand is
$\C[\alpha,\beta,\gamma]/J_g$, where
$J_g=(R_g^1,R_g^2,R_g^3)$.  With
$(R_0^1,R_0^2,R_0^3)=(1,0,0)$, the recurrences of
\cite[Theorem~16]{Munoz1999Ring} are
\begin{align}
 R_g^1&=\alpha R_{g-1}^1+(g-1)^2R_{g-1}^2,
                    \label{eq:caps:munoz-first}\\
 R_g^2&=(\beta+(-1)^g8)R_{g-1}^1
                      +\frac{2(g-1)}gR_{g-1}^3,
                    \label{eq:caps:munoz-second}\\
 R_g^3&=\gamma R_{g-1}^1.\label{eq:caps:munoz-third}
\end{align}
The extreme points of this quotient are
\[
 (a_0,b_0,0)=
 \bigl(\pm4(g-1)i^g,\,(-1)^{g-1}8,\,0\bigr).
\]
Here is an explicit check of the two polynomial facts required for
localization.  Put $\zeta_r=R_r^1$.  Eliminating the other two
polynomials from the recurrence gives
\[
 \zeta_{r+1}=\alpha\zeta_r
 +r^2(\beta+(-1)^r8)\zeta_{r-1}
 +2r(r-1)\gamma\zeta_{r-2}.
\]
For $f_r(\alpha)=\zeta_r(\alpha,b_0,0)$, induction gives
\begin{align*}
 b_0=-8:\quad
 f_{2m}(\alpha)&=\prod_{j=1}^{m}
                   \bigl(\alpha^2-16(2j-1)^2\bigr),&
 f_{2m+1}(\alpha)&=\alpha\prod_{j=1}^{m}
                   \bigl(\alpha^2-16(2j-1)^2\bigr),\\
 b_0=8:\quad
 f_{2m}(\alpha)&=\alpha^2\prod_{j=1}^{m-1}
                   (\alpha^2+64j^2),&
 f_{2m+1}(\alpha)&=\alpha\prod_{j=1}^{m}
                   (\alpha^2+64j^2).
\end{align*}
The second even formula is for $m\ge1$, with $f_0=1$ separately.
For even $g$ use the first row, and for odd $g$ the second.
In either case $f_{g-1}(a_0)\ne0$, whereas $a_0$ is a simple root
of $f_g$.  These are the extreme factors in
\cite[Proposition~20]{Munoz1999Ring}.
Thus $R_{g-1}^1$ is a unit in the localized algebra at this point.
Equation~\eqref{eq:caps:munoz-third} first forces $\gamma=0$.
Since $R_{g-1}^3$ is divisible by $\gamma$,
Equation~\eqref{eq:caps:munoz-second} then forces $\beta=b_0$.
The remaining relation $R_g^1(\alpha,b_0,0)=0$ has a simple root at
$a_0$ by the displayed factors.
After localization it therefore forces $\alpha=a_0$.  The local algebra
is $\C$, not a nonreduced extension of $\C$.

More precisely, the full quotient ring decomposition is
\[
 \bigoplus_{k=0}^{g-1}\Lambda_0^kH^3\otimes
       \C[\alpha,\beta,\gamma]/J_{g-k}.
\]
The terms with $k>0$ are primitive exterior-power summands tensored
with the analogous quotient for a smaller genus.
Their surface eigenvalues have strictly smaller modulus.  They contribute
nothing to an extreme local algebra.  This proves the one-dimensional
\emph{generalized} assertion in the quotient theory.

The full determinant-one group splits into the two eigenspaces of the
degree-four involution $\tau$, its positive eigenspace being the
quotient theory above.  The grading character $U$ associated to an
eighth root of unity satisfies
\[
                  SU=iUS,\qquad PU=-UP,\qquad U\tau=-\tau U.
\]
For operators of cohomological degrees two and four one may take
$U|_{V_j}=e^{-\pi i j/4}$; the opposite grading convention reverses
the character.  It exchanges the two sectors, multiplies the surface
eigenvalue by $i$, and reverses the point eigenvalue.  It carries generalized
spaces isomorphically to generalized spaces.  Applying it to the two
quotient extremes gives all four points in
\eqref{eq:caps:extreme-spectrum}, with the normalization
\eqref{eq:caps:normalization}.  The same ring presentation gives the
stated upper bound for every remaining eigenvalue.
\end{proof}

\begin{lemma}[A nonseparating pencil]\label{lem:caps:pencil}
After blowups at points disjoint from $A$, the tight surfaces, and $Y_0$,
the closure $X_0$ admits a genus-$g$ Lefschetz fibration with a section,
$g\ge2$, all of whose vanishing cycles are nonseparating.  Its positive
monodromy word is the identity in the mapping class group of a surface
with one marked point.
\end{lemma}

\begin{proof}
Rationally approximate and scale the symplectic form, and take a
sufficiently high-degree Donaldson pencil \cite{Donaldson1999}.
Let $F_0$ be its fiber and $b=F_0^2$ its positive number of base points.
Its genus $g_0$ is large.  Since $b^+(X_0)>1$, Taubes's canonical-class
inequality $K_\omega\cdot[\omega]\ge0$
\cite[Theorems~1--2]{Taubes1995} and symplectic adjunction give
\[
       2g_0-2=F_0^2+K_\omega\cdot F_0\ge b.
\]
Use \emph{full} degree doubling: every base point is used in joining the
two copies of the punctured old fiber.  In the doubling monodromy formula,
the new vanishing cycles are nonseparating, and an old separating cycle
remains separating only if one of its sides contains no used base point
\cite[Theorem~4]{AurouxKatzarkov2008}.  If $C$ is a component of a
reducible pencil fiber, then
\[
 F_0\cdot C=k\int_C\omega>0.
\]
A distinct fiber can meet $C$ only at base points.  Both sides of every
old separating cycle therefore contain used base points, so none remains
separating after full doubling.  This is also the mechanism in
\cite[Theorem~3.10]{Smith2001}.

The finite base-point set can be chosen disjoint from the specified
surfaces and hypersurface.  Blow it up.  Any exceptional section marks a
point in every regular fiber; the critical handles are disjoint from this
section.  The monodromy around the sphere is consequently the identity
in the pointed mapping class group.  The resulting fibration is denoted
$\widetilde X_0\to S^2$.
\end{proof}

\begin{lemma}[A model word on an elliptic surface]\label{lem:caps:model}
For every $g\ge2$ there is a genus-$g$ Lefschetz fibration with a section
on a manifold deformation equivalent to $E(g+1)$, all vanishing cycles
nonseparating, whose Donaldson series for the section determinant has a
nonzero coefficient at $e^{(2g-2)t}$ on the fiber line.
\end{lemma}

\begin{proof}
Put $n=g+1$.  Consider the double cover of
$\mathbb P^1_x\times\mathbb P^1_y$ branched in a smooth curve of
bidegree $(4,2n)$, and project to the $x$-line.  A regular fiber is a
double cover of the $y$-line branched at $2n$ points, hence has genus
$n-1=g$.  Fix $y_0$ and impose on the branch polynomial $b$ the condition
\[
                         b(x,y_0)=s(x)^2,
\]
where $s$ is a quadratic section with two simple zeros.  The equations
$z=\pm s(x)$ above $y=y_0$ give sections.  At a zero of $s$, require
$\partial b/\partial y\ne0$.  The total branch curve is then smooth
there, its projection to $x$ is unramified, and each displayed section is
smooth.  The fact that the two sections meet at these points causes no
failure of either section to be a section.

Here is the required genericity away from that slice.  Variations
vanishing on $y=y_0$ are divisible by its defining linear form; their
remaining bidegree is $(4,2n-1)$.  They supply first point jets, second
vertical jets, and the first vertical jets at two distinct points of any
other fixed fiber.  Thus singular branch points impose three independent
conditions on a two-dimensional point space; triple vertical roots
likewise impose three; and two double roots on one fiber impose four
conditions on a three-dimensional space of a fiber and two points.
A generic member of the imposed affine system has none of these
phenomena.  The same argument in the projective coordinate charts covers
points at infinity.  All singular fibers therefore arise from a single
simple branch collision.  The normalization of such a double cover has
$2n-2>0$ remaining branch points and is connected.  Its node is
nonseparating.

The projection to the $y$-line has elliptic regular fiber $F$.
To use the elliptic-surface Donaldson series, we must identify the
surface together with this fiber class.  We first identify a smooth double cover of bidegree
$(4,2)$ with the cubic-pencil $E(1)$, then use cyclic base change and
deformation to obtain the required model of $E(n)$.
Choose a generic pencil
\(y_0A(x,z)+y_1B(x,z)=0\) of curves of bidegree \((2,2)\) on \(\mathbb
P^1_x\times\mathbb P^1_z\). Its eight basepoints are distinct and transverse.
Its incidence surface is the blowup at these basepoints. Projection of the
incidence surface to \(\mathbb P^1_x\times\mathbb P^1_y\) is a finite double
cover. Indeed, at a fixed \(x\) the two quadratic coefficient vectors are
freely variable, and their being proportional is the rank-at-most-one
condition on a \(3\times2\) matrix, of codimension two. Allowing \(x\) to vary
adds only one dimension, so a generic pair has no proportional coefficient
vectors at any \(x\). The three coefficients of the quadratic in \(z\)
therefore have no common zero on the \((x,y)\)-base. Its discriminant has
bidegree \((4,2)\). The finite map is flat, since its smooth surface source is
Cohen–Macaulay and its target is a smooth surface. Locally a flat double cover
is \(w^2=b\); its smooth total space therefore implies that its discriminant
is smooth.

Write \(H_x,H_z\) for the ruling classes and \(E_1,\ldots,E_8\) for the
exceptional classes. The system \(|L|\), where \(L=H_x+H_z-E_1\), is the
resolution of projection of the Segre quadric from the first basepoint to
\(\mathbb P^2\). It contracts the two ruling transforms \(A_1=H_x-E_1\),
\(A_2=H_z-E_1\) and the seven curves \(E_2,\ldots,E_8\). These are nine
disjoint exceptional curves; genericity ensures that none of the other seven
basepoints lies on those two rulings. The elliptic pencil has fiber \[
F=2H_x+2H_z-\sum_{i=1}^{8}E_i   =3L-A_1-A_2-\sum_{i=2}^{8}E_i. \] Thus the
blowdown sends it to a plane cubic pencil with nine simple basepoints. This is
the usual cubic-pencil construction of \(E(1)\), and each blown-up basepoint
gives a section.

Choose a disk in the \(y\)-sphere containing no critical value and choose two
distinct interior points. The connected cyclic degree-\(n\) cover of the
sphere branched at these two points pulls the discriminant back to a smooth
curve of bidegree \((4,2n)\): over the branch values the elliptic fibers are
regular, so the discriminant is transverse to those ruling fibers. Over the
complementary disk the base cover consists of \(n\) disjoint disks. The
inverse image of the chosen disk is a sphere with \(n\) boundary circles,
since its Euler characteristic is \(n-2(n-1)=2-n\). The elliptic fibration
over the original disk is a product. Its pullback is therefore the same
elliptic fiber times this punctured sphere, attached to the \(n\) copies over
the complementary disk with the product identifications. This is the untwisted
iterated fiber sum of \(n\) copies of the cubic-pencil \(E(1)\), namely
\(E(n)\). The base-changed section remains a section, so there are no multiple
fibers.

Finally choose one \(y\)-ruling fiber transverse to both this pullback
discriminant and the smooth discriminant constructed above with the required
genus-\(g\) sections. Smooth divisors of bidegree \((4,2n)\) transverse to
this fixed ruling fiber form a nonempty Zariski-open subset of the projective
linear system, hence are path connected. Over the corresponding connected
open subset of nonzero sections in
$H^0(\mathbb P^1\times\mathbb P^1,\mathcal O(4,2n))$, the equations
$z^2=b$ in the fixed line bundle $\mathcal O(2,n)$ form a proper holomorphic
family of smooth double covers. Thus our cover is complex deformation
equivalent to the cyclic-base-change cubic-pencil model of $E(n)$.
Along a smooth path within this open set the double covers form a proper
smooth family; the inverse images of the
chosen ruling fiber form a smooth subfamily. A lift of the parameter vector
field tangent to this subfamily integrates to an orientation-preserving
diffeomorphism carrying the elliptic fiber to the elliptic fiber. Thus our
model is diffeomorphic to \(E(n)\) with its required class \(F\) preserved.

If $G_f$ denotes the genus-$g$ fiber, then $F\cdot G_f=2$ and
$G_f^2=0$.  The elliptic-surface Donaldson calculation has extreme basic
classes $\pm(n-2)F$ with nonzero coefficients, all other basic classes
having smaller evaluations on $G_f$
\cite{KronheimerMrowka1995}; the elliptic series is computed in
\cite[Section~6, Theorem~6.7, arXiv version~1]{FintushelStern1997},
with its parameters $p=q=1$.  In the convention with
determinant even on $F$, its series has the factor
$\sinh(F)^{n-2}$, times a nonzero constant and $\exp(Q/2)$.
Either genus-$g$ section $z=\pm s(x)$ above $y=y_0$ lies in an
elliptic fiber, so its dual determinant has zero evaluation on a
generic $F$; other determinant
conventions change the nonzero basic coefficients only by signs.
Therefore
\begin{equation}\label{eq:caps:model-series}
 \mathcal D^{\rm section}_{E(n)}(tG_f)
 :=D^{\rm section}_{E(n)}\bigl((1+x_0/2)e^{tG_f}\bigr)
\end{equation}
has nonzero extreme exponential coefficients at
$\pm2(n-2)=\pm(2g-2)$, as claimed.
\end{proof}

\begin{proof}[Proof of Proposition~\ref{prop:caps:nonvanishing}]
Let $g$ be the genus of the pencil on $\widetilde X_0$ supplied by
Lemma~\ref{lem:caps:pencil}.  Remove products over two regular disks
from the genus-$g$ model fibration of
Lemma~\ref{lem:caps:model}.  Its remaining Lefschetz word is an annular
cobordism with both ends identified with the marked product
$\Sigma_g\times S^1$.  Use the determinant dual to its actual section,
including its induced collar identifications.  Let $M$ be its ordinary
map on $V_g$, and let $u_g,v_g$ be the state and costate of the two marked
product-disk pieces.  The annular map commutes with $S$ and $P$:
represent a fiber or point insertion on either product collar and move
it across the annulus.  The even degrees introduce no insertion-order
sign.  Ordinary admissible compactness, the support budget in
Lemma~\ref{lem:caps:labelled-compactness} below, and fixed-transition
gluing of Lemma~\ref{lem:caps:irreducible-gluing} justify this operation.
In particular no simple-type assertion about a relative piece is being
used.

By primary decomposition for two commuting endomorphisms of a
finite-dimensional vector space, their simultaneous generalized spectral
projectors are polynomials in $S,P$.  The nonzero extreme exponential
coefficient in \eqref{eq:caps:model-series}, computed by capping $M$,
therefore gives an extreme line $V_*$ on which
\begin{equation}\label{eq:caps:three-factors}
                    v_g\Pi_*M\Pi_*u_g\ne0.
\end{equation}
Lemma~\ref{lem:caps:extreme} says that this space has dimension one.
Consequently the projected disk state, projected disk costate, and
restriction of $M$ are all nonzero; the last is a nonzero scalar.

The conclusion does not depend on arranging a common sign for all
monomial gluing conventions in advance.  Indeed, consider the two
sequences with $j$ fiber insertions and respectively zero and one point
insertion.  The extreme exponentials in \eqref{eq:caps:model-series}
imply that at least one sequence has upper exponential growth rate at
least $2g-2$.  Multiplying its terms by orientation signs preserves this
rate.  If every extreme projection of the disk--annulus--disk composition
vanished, its terms would instead be bounded by a polynomial in $j$
times $\rho^j$ for some $\rho<2g-2$, by the Jordan decomposition.
This is impossible.  Thus some extreme line always satisfies
\eqref{eq:caps:three-factors}.  The normalization of $S,P$ is fixed by
\eqref{eq:caps:normalization}; the ordinary gluing has no additional
factor accumulating with each insertion.

Let $a_1\cdots a_N$ be the pointed positive word of
$\widetilde X_0$.  Every $a_i$ is a twist on a nonseparating curve away
from the mark.  These curves form one orbit of the pointed mapping class
group.  For each $i$ conjugate a copy of the model word so that a selected
first twist is $a_i$.  A marked oriented fiber identification, a compatible
determinant identification, or continuation induces an isomorphism on
Floer homology commuting with $S,P$.  Its restriction to $V_*$ is
nonzero.  Each conjugated model word thus acts by a nonzero scalar on the
same $V_*$, and so does their product.

Extract the selected twists to the front in their prescribed order.
The elementary Hurwitz exchange for a preceding positive word $R$ is
\begin{equation}\label{eq:caps:hurwitz}
                         Ra=a(a^{-1}Ra).
\end{equation}
It leaves the selected $a$ unchanged and replaces $R$ by a positive
conjugated word.  Repeated use gives the factorization
$(a_1\cdots a_N)R'$ of the product of the model words.  The whole word
and its prefix are pointed identity words; therefore so is $R'$.
Cutting between them gives marked product boundaries, and the composite
of their maps is nonzero on $V_*$.  Both factors, in particular the
desired prefix, are therefore nonzero there.

We spell out why this word operation preserves the data needed by the
gauge theory.  A pointed Lefschetz cobordism is built from a marked product
by attaching the ordered Lefschetz two-handles away from the section
track.  Hurwitz changes are changes of vanishing paths for those handles;
conjugation is a change of marked fiber identification.  The section
track remains straight in the attaching coordinates.  The prefix is
therefore the cobordism of $\widetilde X_0$, with its section, up to its
boundary identifications.  There is no additional loop-of-identifications
invariant: contractibility of $\operatorname{Diff}_0(\Sigma_g)$ and
the point-evaluation fibration, with its fiber based at the identity, give
$\pi_1(\operatorname{Diff}_0(\Sigma_g,p))=\pi_2(\Sigma_g)=0$
\cite[Theorem~1 and Section~1.1, Proposition~1]{Gramain1973}.
Capping the marked section is allowed to have nonzero normal Euler
number.  At every step the determinant is dual to this tracked section,
with its fixed transition functions; we never replace it by a sum over
possible seam lifts.

One can track the normal Euler number explicitly by replacing the marked
point with a small disk.  Each twist away from the mark has its canonical
lift fixing the disk boundary.  A pointed identity word has product
$t_\partial^e$ in this boundary-fixing mapping class group, where
$t_\partial$ is the central boundary twist and $e$ records the section's
normal Euler number (with the convention $e=-[\text{section}]^2$).
The Hurwitz exchange \eqref{eq:caps:hurwitz} holds for these canonical
lifts.  Lifts of a pointed conjugator differ by a central boundary
twist, which cancels in conjugation.  The extracted prefix therefore has
exactly the canonical lifted word, and hence the section Euler number,
of the original pencil.  The complement carries the sum of the model
exponents minus the prefix exponent; no factor is required
to have exponent zero.  Its actual section determinant and boundary
identifications are transported throughout.  A change of marked product
identification transports the disk states by an isomorphism commuting
with $S,P$, so their nonzero projections to $V_*$ remain nonzero.

Cap the prefix with the two nonzero projected disk states.  Since
$V_*$ is a line, the resulting projected composition is nonzero.
The projector is a polynomial in $S,P$, so some monomial of fiber and
point insertions gives a nonzero ordinary Donaldson invariant on
$\widetilde X_0$.  Insertions may be placed on separate product collars.
Their exact geometric representatives and closed interpretation are
justified in Lemmas~\ref{lem:caps:labelled-compactness}
and~\ref{lem:caps:closed-comparison} below; those lemmas apply to these
admissible product boundaries as well as to $Y_0$.

We have obtained a nonzero invariant on $\widetilde X_0$ for its
section determinant.  It remains to pass to the prescribed determinant
$w$ and remove the pencil blowups.  The tight surface $B$ survives these
blowups, so the tight-surface simple-type theorem
\cite[Theorem~8.1]{KronheimerMrowka1995} applies to $\widetilde X_0$, which is simply connected with odd $b^+>1$.
A nonzero invariant then implies a nonzero $(1+x_0/2)$-series.
To see that no cancellation is hidden in this assertion, simple type
reduces every point power using $x_0^2=4$, while, for fixed surface degree
$a$, the terms $D(h^a)$ and $D(x_0h^a)$ have degrees differing by four.
Only one can satisfy the fixed determinant index congruence modulo eight.

The basic-class structure theorem expresses this series as
\begin{equation}\label{eq:caps:structure-series}
 \mathcal D^{w'}(h)=e^{Q(h)/2}
   \sum_j(-1)^{((w')^2+K_j\cdot w')/2}\beta_j e^{K_j\cdot h},
\end{equation}
with distinct $K_j$ and nonzero $\beta_j$ independent of $w'$
\cite[Theorem~1.7]{KronheimerMrowka1995}.
Changing to the pullback of $w$ changes coefficients only by signs and
cannot annihilate this sum of distinct exponentials.  The pullback
blowup formula~\cite[Proposition~1.9]{KronheimerMrowka1995} has a
nonzero exceptional factor, so nonvanishing of
that pullback series implies nonvanishing of the series on $X_0$.
Thus some homogeneous polynomial $h\mapsto D^w_{X_0}(x_0^m h^a)$
is nonzero.

The affine integral lattice $\{h:h\cdot A=1\}$ is nonempty because
$A$ is primitive, and is Zariski dense in its affine hyperplane.  A
nonzero homogeneous polynomial cannot vanish on that whole hyperplane:
rescaling it would make it vanish on the open set $h\cdot A\ne0$.
There is therefore an integral point in the hyperplane at which the
polynomial is nonzero.  For $a=0$ any integral point of that hyperplane
will do.  This proves \eqref{eq:caps:initial-invariant}.
\end{proof}

\subsection{Labelled probes and their exact representatives}

The closure is cut along $Y_0$, while the fixed caps must meet along the
rational homology sphere $Y_2$.  Inserting a positive number of copies of
the bridge $T=f_-\phi^{-1}f_+$ from
Corollary~\ref{cor:ordinary:bridge} creates such a seam.  To retain a
nonzero invariant we must also carry the surfaces across these copies;
the ordinary map $T$ alone does not record their insertions.  We keep
the occurrences separately labelled and construct a finite deformation
of the Floer complex that records them.  Its algebra will show that some
positive power of the deformed bridge still has a nonzero pairing.

Represent the $a$ occurrences of $h$ in
\eqref{eq:caps:initial-invariant} by separately labelled oriented
surfaces.  Their intersections with $Y_0$ all represent the meridian
class in $H_1(Y_0;\Z)$, since $h\cdot A=1$.  A surface cobordism
between homologous one-cycles in $Y_0$, inserted in a collar together
with its reverse, changes each intersection to a prescribed meridian
$\gamma_i$.  Choose the $\gamma_i$ to be disjoint parallel loops.
General position and tubing give pairwise transverse representatives
with no triple intersections.  Put the $m$ point representatives in one
cap, disjoint from all of them.

The two handle halves of $T$ are simply connected, so the meridian
loops bound relative surfaces on the two halves separately.  Fix these
surfaces, joining the corresponding labels at the two ends of $T$.
Choose its determinant with evaluation one on both end copies of $A$ and the prescribed trivialization at its internal
rational seam.  The increasing-handle coordinates in
Section~\ref{sec:ordinary} give precisely these restrictions.  Glue it
to $w$ and repeat the same determinant transitions for every copy of
$T$.  The closed manifold obtained by inserting $T^k$ in $X_0$ at
$Y_0$ is denoted $X'_k$, with $X'_0=X_0$.

\begin{lemma}[Exact cuts and the low-dimensional support budget]
\label{lem:caps:labelled-compactness}
For the labelled surfaces just described, and also for the separated
fiber and point insertions used above, the exact adjoint tests can be
chosen coherently on cylinders, caps, and cobordisms.  In cut dimensions
zero and one, their compactifications have only the ordinary irreducible
breaking strata.  No curvature particle contributes.  These statements
hold for the parameterized comparisons of exact tests as well.
\end{lemma}

\begin{proof}
At a fixed or moving point, use the dependence locus of two sections of
the universal complexified adjoint bundle of rank three.  Its principal
stratum has complex codimension two, hence real codimension four, and
represents
\[
                     c_2(\ad E\otimes\C)=-p_1(\ad E).
\]
The zero-rank stratum has larger codimension and is avoided in the
dimensions at issue.  A moving point on an oriented surface contributes
two position parameters, so its net cut degree is two.  A fixed point
has degree four.

The section construction of Proposition~\ref{prop:analysis:ordinary-library}
uses holonomy and frame-equivariant sampled transports from the marked
point to finite systems of loops detecting irreducibility.  The center
acts trivially on the adjoint output.  Irreducibility gives full value
variation at the output, hence transverse choices of the two sections.
On a cylinder the choices are translation invariant relative to the
mark; on cobordisms they approach those choices at the ends.  The
sampling terms carry small-ball energy cutoffs along their paths.
For thin paths one uses arbitrarily fine path libraries, with sufficiently
rapidly decreasing coefficients.  The cutoff limits include curvature
atoms in the limiting energy measure.  Thus a term either has its
prescribed convergent off-atom transport or is cut off when its sampling
neighborhood acquires an atom.  Loop systems avoiding a given finite
set of atoms still detect irreducibility and still give independent
value variations.  The resulting exact cut condition persists unless
its marked point itself hits a curvature atom.  On a particle stratum
the retained equation is allowed to depend on its retained atom
positions.  Indeed a smooth sampling cutoff replaces its energy
argument by
\[
 \int\rho\,|F_a^0|^2+8\pi^2\sum_j\ell_j\rho(p_j),
\]
which is smooth in the positions $p_j$ on that stratum.  Thus persistence
means convergence to a section on the residual-field-and-particle
stratum, not necessarily to the ordinary section on the residual field
alone.  The latter equality is not needed for the dimension budget.
This is the
support-exact feature for which the construction, rather than a
representative depending indiscriminately on a whole surface
neighborhood, is needed here.

Choose parameters transverse simultaneously on all lower-charge strata,
including remaining marks and the positions of particles.  This remains
possible for the particle-dependent sections just described: at any
such configuration a detecting graph and sufficiently small sampling
balls can avoid the finite atom set, and the associated surviving term
has unrestricted adjoint output variation.  The particle positions are
counted once among the stratum parameters, not again as new cut data.
Ends carrying
the odd determinant on $A$ are admissible.  Every nonconstant cylinder
level and every cobordism piece considered in this labelled construction
is consequently irreducible.  A constant critical cylinder cannot carry
a generic mark: modulo translation its mark has only the one loop
coordinate, against a codimension-four condition.  Floer compactness
therefore assigns any escaping mark to a nonconstant trajectory level,
unless that mark hits an escaping particle.

For a compact-part particle, the precise budget is as follows.  The
``Released degree'' column records the net degree of discarded cuts, after their
mark parameters are included.
\begin{center}
\begin{tabular}{lrrr}
\toprule
Position of the particle & Released degree & Position dimension & Sum\\
\midrule
Outside the supports & $0$ & $4$ & $4$\\
On one surface & $2$ & $2$ & $4$\\
At two transverse surfaces & $4$ & $0$ & $4$\\
At a fixed point & $4$ & $0$ & $4$\\
\bottomrule
\end{tabular}
\end{center}
There are no triple intersections or point--surface incidences.  If the
lost charge is $\ell$ and it is carried by $r$ distinct particles, then
$r\le\ell$, while the ASD index drops by $8\ell$.  The residual cut
stratum, with particle positions retained, therefore has dimension at
most
\begin{equation}\label{eq:caps:bubble-budget}
                d-8\ell+4r\le d-4\ell,
\end{equation}
where $d$ is the original cut dimension, including any comparison
parameter.  This is negative for $d\le1$ and $\ell>0$.
The same reasoning applies on extracted nonconstant cylinder levels.
If an atom escapes along a constant-flat region with no nonconstant
level retained there, it can discard at most one label, since the
parallel cylindrical supports $\R\times\gamma_i$ are disjoint.
All equations on retained levels persist, so this case is at least as
strongly excluded.  This is the incidence estimate of
Lemma~\ref{lem:analysis:incidence} in the present setting.

At the remaining broken configurations, exponential convergence to
nondegenerate critical points and
Lemma~\ref{lem:caps:irreducible-gluing} give exactly
the products of the cut factors.  The increasingly distant sampling
offsets may be truncated on long gluing collars with faster than
exponential tails; their errors then lie within the ordinary gluing
estimates.  Section choices and determinant-line orientations are
compatible with these products.  Thus the only ends of an oriented
one-dimensional cut moduli space are its ordinary irreducible breaks
and its prescribed comparison endpoints.
\end{proof}

\begin{lemma}[Comparison with closed Donaldson insertions]
\label{lem:caps:closed-comparison}
The closed counts obtained from these exact tests and the fixed
transition functions are the Donaldson invariants for that determinant,
with each surface and point class replaced by its multiple $4\mu$.
They have no additional sum over determinant lifts.  The assertion
applies to $X'_k$ and to the fiber gluings used in
Proposition~\ref{prop:caps:nonvanishing}.
\end{lemma}

\begin{proof}
All the $X'_k$ are simply connected: the original two halves and the
handle halves of $T$ are simply connected, and Van Kampen applies at
each connected seam.  They retain a tight surface and $b^+>1$.
Their determinants remain odd on the retained copies of $A$.
The same properties needed here hold for the simply connected model and
pencil manifolds.  We first explain comparison from long admissible
collars, where the Floer perturbation was chosen before the gluing
length.

A reducible would have integral line difference $d$ with odd evaluation
on every odd-class slice.  On a fixed slice metric its tangential
curvature therefore has a uniform positive $L^2$ lower bound, by pairing
with a fixed dual form and Cauchy--Schwarz.  Integrated along a long
product collar this lower bound grows linearly with its length.  The
near-ASD equation, with gradient-type Floer perturbation and sufficiently
small remaining error, converts it to a growing lower bound for the
charge on that product.  Other admissible collars have the same
property, and their gradient energy identities have endpoint terms
that telescope.  In a fixed total charge range this rules out a
reducible while the Floer perturbations are turned off on collars that
are still sufficiently long.  It is therefore legitimate to perform
this first step with the earlier fixed small Floer data; no reversal of
the order of smallness and gluing length is required.  Thereafter choose
a generic path to an unperturbed closed metric, with perturbations as
small as needed.  Generic periods exclude the finitely relevant
reducible classes throughout this path; $b^+>1$ permits a generic
one-parameter path.  The bubble budget
\eqref{eq:caps:bubble-budget} applies to these comparison homotopies.

Localize the exact section data near the surface representatives and in
small disjoint point balls at the final generic unperturbed metric.
An irreducible ASD connection on one of these simply connected
manifolds is locally irreducible: a local parallel line extends by ASD
unique continuation, and simple connectivity removes monodromy of this
extension.  An irreducible flat exception cannot occur.  Thus the
localized choices retain the needed independent transversality.
For a fixed point the adjoint rank-drop locus is already the usual
geometric representative of $-p_1$ on its ball.

For a surface compare one test at a time.  Omit that degree-two test;
the remaining partially cut ASD moduli space has dimension two.
It is compact.  Indeed, take the chosen small neighborhoods so that
at most two surface neighborhoods meet at any point and no point ball
meets a surface neighborhood.  If a particle occurs, discard the tests
whose neighborhoods contain it.  At most degree four is released per
particle.  For this argument project to the residual ASD solution and
ignore particle-position parameters: every sampling graph, output,
and cutoff for a retained test lies in its atom-free neighborhood.
Consequently these retained equations are exactly the ordinary
lower-level local equations and are independent of particles outside
their neighborhoods.  This localization, unlike the earlier general
particle-stratum argument, is what permits forgetting the positions.
The residual problem has dimension at most
$2-8\ell+4\ell<0$.  Hence it has no solutions.  This establishes the
compactness needed for the comparison, rather than assuming compactness
of an arbitrary partially cut moduli space.

On this compact two-dimensional space the moving-point intersection
number is exactly the slant product of $c_2(\ad E\otimes\C)$ with
the surface.  By \eqref{eq:caps:normalization} it is $4\mu$ of that
surface.  Replacing the surface tests one at a time identifies the full
count with the ordinary Donaldson intersection number, and the point
tests contribute $4\mu(x_0)$ in the same normalization.  Generic
homotopies of these exact tests have no extra ends by
Lemma~\ref{lem:caps:labelled-compactness}.

Finally, the glued transition functions are fixed actual
determinant-preserving bundle identifications, not arbitrary projective
identifications with a prescribed characteristic class.  A projective
seam automorphism has an obstruction in $H^1(Y;\Z/2)$ to lifting to
the determinant-one gauge group.  A nonzero obstruction excludes that
automorphism from the matching group at the outset, even if it would
preserve $c_1$ and $c_2$.  The full determinant-one Floer complex likewise
keeps distinct the orbits related by such a nonliftable twist.

For the remaining, liftable automorphisms, obstruction theory on the
three-manifold identifies components of the determinant-one gauge group
with $H^3(Y;\Z)$: its fiber $\SU(2)$ is two-connected, and conjugation
acts trivially on $\pi_3(\SU(2))=\Z$.  Thus a degree-zero allowed
automorphism is homotopic to the identity.  A nonzero degree changes
$c_2$ and changes the ASD index by eight times that degree.  The actual
zero-dimensional index fixes this remaining integer.  In the pointed
word construction the actual section and its dual determinant transport
these identifications through every conjugation and Hurwitz move.
Therefore ordinary gluing here counts the required fixed determinant
problem; neither a hidden $H^1(Y;\Z/2)$ twist nor a new sum of lifts is
introduced.

This restriction on matching isomorphisms does not omit closed
connections.  Fix the actual closed rank-two bundle $E$ with the
prescribed determinant and second Chern class, and restrict it to the
pieces using the transition induced by $E$.  Every closed connection on
$E$ has literal restrictions with equal full determinant-one boundary
orbits.  Compactness and gluing are performed in that gauge convention.
A projective change of gauge with nonzero lifting obstruction can move
one full boundary orbit to another; the composition sums over all full
orbits and does not quotient by this action.  Thus such boundary states
are retained, rather than their corresponding closed solutions being
discarded.  Ordinary maps also retain all compatible relative bundle
sectors.  Their integer degree bookkeeping sums to $c_2(E)$ at the
selected total index.  Conversely a full-state match with these fixed
determinant transitions and this total degree glues a rank-two bundle
with $c_1(E),c_2(E)$, hence the prescribed closed bundle, since these
classes classify rank-two complex bundles over a four-dimensional
complex.  The usual determinant-one gauge-matching quotient identifies
its connections with the closed gauge orbits, and
Lemma~\ref{lem:caps:irreducible-gluing} gives multiplicity one.  This
applies in particular to the annulus--disk compositions above.

For comparison, the published
Fukaya--Floer gluing formula in \cite[Theorem~3.1]{Munoz1999FukayaFloer}
uses a convention permitting an additional nonliftable seam twist and
has a corresponding sum of two determinant invariants.  We have instead
constructed and glued the determinant-one counts with fixed collars;
that published sum is not being invoked as a fixed-lift formula.
\end{proof}

\subsection{The finite deformation and a positive power of the bridge}

Let $C$ be the ordinary integral admissible Floer complex of $Y_0$,
with its fixed determinant.  Tensor with $\Q$ when needed.  Keep the
labels distinct and set
\begin{equation}\label{eq:caps:label-ring}
 L_a=\Q[t_1,\ldots,t_a]/(t_1^2,\ldots,t_a^2),
 \qquad J=(t_1,\ldots,t_a).
\end{equation}
The variables have even parity (one may give each degree $-2$).
For a subset $P\subset\{1,\ldots,a\}$ write $t_P=\prod_{i\in P}t_i$.
For $a=0$, use $L_0=\Q$.

\begin{proposition}[Finite labelled Floer construction]\label{prop:caps:labelled}
There are a differential $d_L$ on $C_L=C\otimes L_a$, a cycle
$u\in C_L$, a cocycle $v\in\Hom_{L_a}(C_L,L_a)$, and a chain map
$T_L:C_L\to C_L$ such that $T_L\bmod J=T$.  The all-label coefficient
of $vT_L^k u$ is, up to a fixed nonzero normalization and orientation
sign, the closed Donaldson invariant on $X'_k$ with the indicated
labelled closed surfaces and fixed point insertions.  In particular,
\begin{equation}\label{eq:caps:coefficient-zero}
                    [t_1\cdots t_a](vu)\ne0.
\end{equation}
\end{proposition}

\begin{proof}
For each subset $P$, count cylinder trajectories with one moving point
on $\R\times\gamma_i$ and its codimension-four adjoint test for every
$i\in P$.  A rigid such trajectory has unparametrized index
$2|P|$, or parametrized trajectory index $1+2|P|$.
Let its count be $d_P$, and set
\[
                         d_L=\sum_Pt_Pd_P.
\]
Do not retain a mark or a test for a label outside $P$.  Define the cap
coefficients $u_P,v_P$ by the corresponding relative surfaces in the
caps and their fixed point insertions.  Define $(T_L)_P$ by the
relative surfaces in $T$.  All these counts use the exact tests of
Lemma~\ref{lem:caps:labelled-compactness}.

The oriented boundary of a one-dimensional cylinder space is a broken
pair with each label assigned to exactly one factor.  Its algebraic
identity is
\[
              \sum_{P_1\sqcup P_2=P}d_{P_1}d_{P_2}=0.
\]
This is exactly the coefficient of $t_P$ in $d_L^2$, since repeated
labels vanish in $L_a$.  The same boundary calculation on a cap or a
cobordism gives $d_Lu=0$, $vd_L=0$, and
$d_LT_L=T_Ld_L$, with the prescribed common grading sign if that
convention is used.  End-generator orientation lines are paired by
ordinary determinant-one composition.  All insertion degrees are even,
so distributing labels introduces no additional Koszul sign.  There
are no factorials: these are distinct labels, each used at most once.

Stretch the $Y_0$ seams in the closed manifold.  Its ends are admissible.
Lemma~\ref{lem:caps:labelled-compactness} excludes particles and unwanted
constant levels, and Lemma~\ref{lem:caps:irreducible-gluing} gives the product of the
cut factors, with labels distributed among them.  This is precisely
the coefficient expansion of $vT_L^ku$.  The fixed transition data and
Lemma~\ref{lem:caps:closed-comparison} identify each coefficient with
its actual fixed-determinant Donaldson invariant.  At $k=0$, the
all-label coefficient is the nonzero invariant
\eqref{eq:caps:initial-invariant}, multiplied by $4^{a+m}$ and the
fixed orientation sign.  This proves \eqref{eq:caps:coefficient-zero}.
\end{proof}

\begin{lemma}[Nilpotent deformation and positive powers]\label{lem:caps:positive-power}
There is $k\ge1$ for which
\begin{equation}\label{eq:caps:positive-coefficient}
                    [t_1\cdots t_a](vT_L^ku)\ne0.
\end{equation}
\end{lemma}

\begin{proof}
By Corollary~\ref{cor:ordinary:bridge}, $T$ is a rational
quasi-isomorphism.  Filter the mapping cone of $T_L$ by powers of $J$.
Since $J^{a+1}=0$, this is a finite filtration.  Its associated graded
complex is a direct sum of copies of the rational mapping cone of $T$,
tensored with the finite-dimensional spaces $J^r/J^{r+1}$, and is
acyclic.  Induction through the short exact sequences of filtration
steps shows that the cone of $T_L$ is acyclic.  Therefore the induced
endomorphism $F$ of $H_*(C_L)$ is an automorphism.  This homology is
finite-dimensional over $\Q$; no freeness over $L_a$ has been asserted
or needed.

Let $N=\dim_\Q H_*(C_L)$ and write the characteristic polynomial of
$F$ as $a_0+a_1s+\cdots+a_Ns^N$.  Its constant coefficient is nonzero
because $F$ is invertible.  Cayley--Hamilton gives
\begin{equation}\label{eq:caps:cayley-hamilton}
               \id=-a_0^{-1}\sum_{j=1}^Na_jF^j.
\end{equation}
Apply the functional $[t_1\cdots t_a]\,v(\,\cdot\,)u$ to this
identity on homology.  Its value on the identity is nonzero by
\eqref{eq:caps:coefficient-zero}.  At least one positive-power value
must therefore be nonzero, proving \eqref{eq:caps:positive-coefficient}.
\end{proof}

\subsection{Moving all probes into the final caps}

Fix a positive $k$ supplied by Lemma~\ref{lem:caps:positive-power}.
The associated nonzero closed count on $X'_k$ has only even-degree
insertions.  The dimension congruence consequently forces $b^+(X'_k)$
to be odd; $b_1=0$ and the ASD dimension is
$8\kappa-3(1+b^+)$.  Split $X'_k$ at an internal rational $Y_2$ seam.
Since $H_1(Y_2;\Q)=H_2(Y_2;\Q)=0$, Mayer--Vietoris gives
\begin{equation}\label{eq:caps:homology-split}
 H_2(X'_k;\Q)=H_2(C_-^0;\Q)\oplus H_2(C_+^0;\Q).
\end{equation}
Expand each probe class with this direct sum and use multilinearity of
the closed invariant.  The resulting finite sum is nonzero, so at least
one allocation of all the probes to the two cap interiors has nonzero
value.  Clear its rational denominators.  Integral surface classes in a
cap have embedded oriented representatives in its interior, after
resolving double points if necessary.  Take separate representatives for
separate labels, arrange pairwise transversality without triple points,
and keep them disjoint from the point probes.

Each $C_\pm^0$ is simply connected by the same handle-half and Van Kampen
argument used above.  It contains its old tight surface $B$, so it has
$b^+>0$, and contains a copy of $A$ on which its determinant evaluates
as one.  Its determinant is trivialized at the rational boundary.
Blow up one further point in the interior of each cap, away from all
these surfaces and probes.  Extend the determinant with zero exceptional
evaluation.  The pullback blowup formula
\cite[Proposition~1.9]{KronheimerMrowka1995}, restricted to the old homology,
preserves the invariant, because its exceptional factor has constant
term one.  Denote the resulting caps by $C_-,C_+$.

\begin{proposition}[Cap periods and the ordinary pairing]\label{prop:caps:periods}
Generic period choices on the old cap parts exclude reducible ASD cap
solutions in the bounded charge ranges needed for the cap cycles,
comparison homotopies, and the closed pairing.  After sufficiently long
finite isolation of the separate exceptional summands, and sufficiently
small ordinary Floer data, stretching the common $Y_2$ seam computes
the nonzero allocated invariant as the pairing of a cycle and a cocycle
on the irreducible complex.
\end{proposition}

\begin{proof}
For a reducible with determinant $c$, its integral line difference is
$d=2l-c$ for some integral $l$.  In particular
\[
                   d(A)\equiv c(A)=1\pmod2,
\]
so $d$ is nonzero in rational cohomology.  A reducible ASD solution with
rational flat end would require $d^+=0$.  Since $b^+(C_\pm)>0$,
generic cap periods avoid this condition.  The precise cap version of
the period argument, including compactly supported metric variation,
small perturbations, bounded charge, and the long isolated exceptional
summand, is Proposition~\ref{prop:estimates:cap}.  Its metric choices
apply here because the old cap contains both $A$ and $B$ and the probes
are supported away from the isolated summand.

One must use its bounded-charge conclusion, rather than claiming that
arbitrary cap data give the same cap states.  Here are the charge bounds
that supply its hypothesis.  Let $z_W$ be the local probe degree on one
cap.  Its cycle entries have actual index $i=z_W$, and the
one-dimensional spaces giving the cycle identity have $i=z_W+1$.
For a comparison with $t\in\{0,1\}$ parameters the corresponding
condition is $i+t-z_W\in\{0,1\}$.  Only these finitely many integer
indices occur.  With fixed reference lifts for the finitely many end
generators $\alpha$, the APS formula takes the form
\[
                       i=8\kappa+b_W(\alpha).
\]
The boundary correction $b_W(\alpha)$ is uniformly bounded as the
ordinary perturbations become sufficiently small.  Indeed the flat
moduli space is compact; elliptic spectral estimates bound uniformly
the number of eigenvalues near zero, so only uniformly finitely many
such eigenvalues can contribute integer spectral-flow jumps under a
small perturbation.  The smooth boundary and transgression terms are
also bounded on this compact family near the flat moduli space.
At a central flat the correction
is explicitly $-3(1+b^+(W))$ by
Equation~\eqref{eq:ordinary:central-index}.  Thus all these cap problems
have a bounded charge range, including reducible candidates before
period exclusion.  The closed pairing has the fixed total charge
\[
                    8\kappa_{\rm closed}
                        =z+3(1+b^+(X)).
\]

These bounds do not deteriorate with a Floer collar length.  On a
gradient collar the perturbed Chern--Simons identity bounds the
topological charge below by the difference of the bounded perturbation
at its two endpoints, with a fixed error from the permitted integrable
non-gradient terms.  Summing through successive segments or broken
levels telescopes those endpoint terms.  Thus the charge of an entire
outgoing chain has a lower bound independent of its length and of the
number of levels.  Subtracting this bound from the fixed total charge
bounds the charge on the compact cap part from above.  The same argument
applies to either cap in the closed pairing, using the fixed local
geometry on the other compact part.  Chern--Weil then bounds the
curvature on compact parts.  On a prefix of length $N$, the remaining
small self-dual error contributes at most a fixed constant times
$\delta^2N$, so the bound has the form
\[
        \int_{\text{cap and prefix through }N}|F_a^0|^2
                     \le E_{\mathcal P}+C_W\delta^2N.
\]
Here $E_{\mathcal P}$ can depend on the fixed finite problem, while
the coefficient of $N$ is controlled by the fixed collar data.  This
is the prefix bound required for the exponentially decreasing harmonic
cutoff tails in Proposition~\ref{prop:estimates:cap}.  It also covers
central outgoing chains; their later dimension exclusion is not used
to obtain this energy bound.  A formal infinite tail may be asymptotic
to a perturbed nonflat connection, so no finite unperturbed curvature
norm on that entire tail is asserted.  Isolation depths and smallness
thresholds can now be chosen for all the fixed index ranges just listed.

Count regular irreducible cap solutions with their local cuts, obtaining
$x$ and $y$.  Bubbling is absent by the exact support budget of
Lemma~\ref{lem:caps:labelled-compactness}.  Reducible cap pieces are
absent by the period choice.  At a central seam the three-dimensional
central stabilizer contributes three matching indices, as in
Lemma~\ref{lem:ordinary:central-cost}.  Each of the two irreducible cap
pieces has nonnegative residual dimension after its local cuts.  Such
a central matching therefore cannot occur in a zero-dimensional closed
count.  More generally a maximal central cylinder group includes its
exterior matchings and has the same surplus of at least three.  At a
central end in a cap-cycle identity it must also have the requisite
nonconstant outgoing trajectory.  It cannot be an end of a
one-dimensional cut cap moduli space.  The only remaining boundary
points are ordinary irreducible Floer breaks.  Their signed sum gives
$\partial x=0$ and $y\partial=0$.

At a sufficiently long finite common neck, every rigid closed solution
is therefore an irreducible match of two rigid cap solutions.  Conversely
each such match glues by Lemma~\ref{lem:caps:irreducible-gluing}.
Fixed determinant transitions identify the glued bundle, and cancelling
adjacent orientation lines gives the closed count up to its fixed
comparison sign.  There is no extra factor of two from the common
central subgroup $\{\pm1\}$ of an irreducible stabilizer: it is already
included in the determinant-one gluing convention.  We obtain the
pairing $yx$, which equals the nonzero allocated invariant up to the
nonzero $4\mu$ normalization.
\end{proof}

\begin{proof}[Proof of Theorem~\ref{thm:caps:pairing}]
Use the caps and probes just constructed and
Proposition~\ref{prop:caps:periods}.  Their topology gives (i) and (ii),
and their ordinary pairing gives (iii).  All exact point and surface
classes are integral multiples of the Donaldson classes.  The remaining
rational factors come from the fixed denominators in
\eqref{eq:caps:homology-split} and, if used, rational perturbation
weights.  Clearing these fixed denominators yields integral cycles and
cocycles in the finite integral complex, while multiplying $yx$ by a
nonzero integer.  Let $z$ be the sum of the real degrees of the resulting
fixed probes.  This proves (iv), with the extension of the cap period
choices supplied by Proposition~\ref{prop:estimates:cap}.

For later metric choices, recompute the pairing by stretching the same
seam in the old closed, blown-up manifold.  Its fixed-determinant closed
Donaldson invariant is unchanged because $b^+>1$; the exact-test
comparison is Lemma~\ref{lem:caps:closed-comparison}.  The cap period
exclusion and central matching budget apply again at those choices.
This proves the last assertion without requiring separate
metric-independence of either individual cap state.
\end{proof}
